% Document A (readable edition). Lean counterpart: lean/GL2TF/VolII/Kernel.lean
% Chapter files of A and B have identical names; owner per sprint: see the sprint note.
% §1 written (M4, unit U9; docs/sprints/M4.md §1.2 items 58--60); §§2--3 (M7a, unit U5; docs/sprints/M7a-plan.md §1.2 items 36--47; stretch item 48 moved to M9b, decision D-LATE;
% Lean: lean/GL2TF/VolII/KernelClasses.lean); §§4 ff. (spectral half): G0 outline pending (M10a).
% The stale G0 outline line at the end of §1 was removed in M7a W1 (lead decision; re-opens II.kernel.thm:integral-operator for G2).
% Kernel normalization: docs/notes/M4-normalizations.md §5 and the notation table.
% Fork algebraic_jl: only §2 (geometric classes) is kept, without the class kernels; title and scope rewritten (../../EDITS.md).

\chapter{Geometric conjugacy classes}\label{ch:ch-03-kernel}

\noindent\emph{Scope:} area \texttt{kernel} (its algebraic part). \emph{Sources:} \cite[§10]{AIntro}; \cite[Lecture II §1]{Gel}.
In the parent project this chapter is about the kernel of the right regular representation and its decompositions; only the
classification of the elements of $\GL_2(\Q)$ by their characteristic polynomials is kept here.

\section{Geometric classes}\label{sec:kernel-classes}

\noindent\emph{Sources for this section:} \cite[§10, pp.~53--54]{AIntro} (the equivalence classes $\mathfrak o\in\mathcal O$; for $\GL(n)$, $\gamma$ and $\gamma'$ are
equivalent iff they have the same characteristic polynomial); \cite[Lecture II §1]{Gel} (the classes for $\GL_2$: central-unipotent, hyperbolic, elliptic).

\begin{definition}[geometric classes]
\tfstmt{II.kernel.def:geometric-class}
\uses{I.conv.def:gl2, I.topgrp.def:gl2-adelic}
\end{definition}
\noindent\emph{Comparison.} This is Arthur's $\mathcal O$ for $\GL_2$ \cite[§10, p.~54]{AIntro}: $\gamma,\gamma'$ are $\mathcal O$-equivalent iff their semisimple parts are
conjugate, iff they have the same characteristic polynomial. Every $(t,d)$ occurs (the companion matrix $\begin{psmallmatrix}0&1\\-d&t\end{psmallmatrix}$ has
class $(t,d)$), and conjugate elements have the same class. A class is a union of conjugacy classes, by \cref{II.kernel.thm:semisimple-class-conjugacy} of one
or two.

\begin{definition}[types of classes]
\tfstmt{II.kernel.def:class-types}
\uses{II.kernel.def:geometric-class}
\end{definition}
\noindent In terms of the discriminant $t^2-4d$: zero, a non-zero square, a non-square in $\Q$. \emph{Examples.} $(2,1)=\mathcal O_1$ contains $1$ and
$n(1)$; $(3,2)=\mathcal O_{2,1}$ contains $\diag(2,1)$; $(0,1)$ ($P=X^2+1$) and $(0,-2)$ ($P=X^2-2$) are elliptic, containing $\gamma_i$ and $\gamma_{\sqrt2}$ of
\cref{ch:ch-04-orbital}. Arthur calls the elliptic classes of $\GL_2$ \emph{anisotropic} \cite[§10, p.~53]{AIntro}.

\begin{theorem}[trichotomy]
\tfstmt{II.kernel.thm:class-trichotomy}
\uses{II.kernel.def:class-types, II.orb.def:elliptic}
\end{theorem}
\begin{proof}
Let $P=X^2-tX+d$ with $d\neq0$. If $P$ has a root $a\in\Q$, then $P=(X-a)(X-b)$ with $b=t-a$ and $ab=d\neq0$, so $a,b\in\Q^\times$, and $\mathcal O$ is
central-unipotent ($a=b$) or hyperbolic ($a\neq b$). Otherwise $P$, of degree $2$ without a root, is irreducible. The types exclude each other: (i) and
(ii) have a root and (iii) has none, and the factorization of $P$ into monic linear factors is unique ($\Q[X]$ is factorial), so $\{z,z\}\neq\{a,b\}$ when
$a\neq b$. The second claim holds because $P_\gamma=X^2-\tr(\gamma)X+\det\gamma$ is the polynomial of the class of $\gamma$, and $\gamma$ is elliptic iff $P_\gamma$ is
irreducible (\texttt{II.orb.def:elliptic} with $K=\Q$).
\end{proof}

\begin{theorem}[conjugacy within a class]
\tfstmt{II.kernel.thm:semisimple-class-conjugacy}
\uses{II.kernel.def:geometric-class, II.kernel.def:class-types, I.conv.def:torus, I.conv.def:borel, I.conv.def:unipotent}
\end{theorem}
\begin{proof}
\uses{II.kernel.thm:class-trichotomy, II.orb.thm:elliptic-iff-not-borel-conjugate}
\emph{Companion form.} If $\gamma\in\GL_2(\Q)$ is not central, then $\gamma$ is conjugate to $C(t,d)=\begin{psmallmatrix}0&1\\-d&t\end{psmallmatrix}$, $(t,d)$ its class. Indeed, as
in the proof of \texttt{II.orb.thm:centralizer-explicit}, some row vector $v$ has $v,v\gamma$ linearly independent (if every $v\neq0$ were an eigenvector, then $e_1$,
$e_2$, $e_1+e_2$ would be, forcing equal eigenvalues and $\gamma$ central), and the matrix $\delta$ with rows $v,v\gamma$ satisfies $\delta\gamma\delta^{-1}=C(t,d)$ by
Cayley--Hamilton ($v\gamma^2=tv\gamma-dv$). Conversely, conjugate elements have the same class. A central element $z\cdot1$ has class $(2z,z^2)$, of type (i).

(a) Elements of a hyperbolic or elliptic class are not central, so all are conjugate to $C(t,d)$, hence to each other. (b) $\diag(p,q)$ has polynomial
$(X-p)(X-q)$, which equals $(X-a)(X-b)$ iff $\{p,q\}=\{a,b\}$. (c) An element of $B(\Q)$ has the root $\gamma_{11}$ of its characteristic polynomial; so it is not
elliptic (\cref{II.kernel.thm:class-trichotomy}; equivalently \texttt{II.orb.thm:elliptic-iff-not-borel-conjugate}). (d) Let $\gamma\in\mathcal O_z$. If $\gamma=z'\cdot1$ is
central, its class $(2z',z'^2)=(2z,z^2)$ gives $z'=z$. Otherwise $\gamma$ is conjugate to $C(2z,z^2)$, and so is the non-central element
$z\,n(1)=\begin{psmallmatrix}z&z\\0&z\end{psmallmatrix}$ of the same class; hence $\gamma$ is a conjugate of $z\,n(1)$. Conversely $z\cdot1$ and the conjugates of $z\,n(1)$ lie in
$\mathcal O_z$.
\end{proof}
\noindent So a hyperbolic class meets $T(\Q)$ (\cref{II.kernel.thm:semisimple-class-conjugacy}(b)) and every hyperbolic class is represented by some $\diag(a,b)$,
$a\neq b$. $\mathcal O_z$ consists of two conjugacy classes, which is
why the class of $1$ contains the unipotent elements.

